\chapter{Pourquoi l'apprentissage automatique en finance est différent}\label{ch:ml:why-financial-machine-learning-is-different}

Un nouveau modèle prévoit le rendement du mois prochain de 500 actions. Hors échantillon, sur dix ans, son R² vaut 0.57~\%, et
le desk est satisfait. Il devrait l'être~: sur le marché simulé sur lequel le modèle a été entraîné, où les rendements espérés sont
connus parce qu'ils ont été introduits artificiellement, la meilleure prévision possible atteint 0.60~\%. Les mêmes arbres boostés par gradient,
appliqués à une tâche dont le signal est fort, expliquent 92~\% de la variance. Rien n'est faux dans l'un ou l'autre chiffre. Le modèle
a fait aussi bien sur les actions que sur la tâche facile~; les actions ne contiennent tout simplement presque rien à apprendre, et ce qu'elles
contiennent change, est partagé par toutes les actions à la fois, et se dissipe dès que le premier venu l'exploite. Ce chapitre fixe
le vocabulaire de l'apprentissage automatique que le reste du livre utilise, et mesure, sur des données dont la vérité est connue, les quatre
façons dont les données de marché diffèrent de celles sur lesquelles le domaine s'est construit.

\section{Apprendre à partir des données~: le vocabulaire}

\begin{definition}[Apprentissage supervisé et non supervisé]\label{def:ml:why-financial-machine-learning-is-different:learning}
L'\emph{apprentissage supervisé}\index{apprentissage supervisé} (supervised learning) ajuste une fonction $f_\theta$ des variables $x\in\mathbb R^p$ vers
une cible $y$ sur des exemples $(x_i, y_i)$, $i = 1,\dots,n$, en minimisant une perte empirique
$\mathcal L_n(\theta) = \frac1n\sum_i\ell(y_i, f_\theta(x_i))$, de sorte que $f_\theta$ prédit $y$ pour un nouveau $x$.
L'\emph{apprentissage non supervisé}\index{apprentissage non supervisé} (unsupervised learning) n'a pas de cible~: il décrit la distribution de $x$ elle-même
(regroupements, facteurs, représentation en dimension réduite, densité).
\end{definition}

La cible est la cible de prévision du livre~7 (chapitre~6)~: un rendement à terme sur un horizon de prévision, ou une étiquette construite à partir
d'un rendement (chapitre~2 de ce livre). La perte $\ell(y,\hat y)$ prend toujours deux
arguments ici~; $\mathcal L_n$ est toujours indicée. Le \emph{taux d'apprentissage} (learning rate) des chapitres suivants s'écrit $\eta$,
un symbole propre à ce livre (la série réserve ailleurs $\eta$ au vol-of-vol).

\begin{definition}[Ensembles d'entraînement, de validation et de test~; hyperparamètre]\label{def:ml:why-financial-machine-learning-is-different:sets}
L'\emph{ensemble d'entraînement}\index{ensemble d'entraînement} (training set) est constitué des données sur lesquelles $\theta$ est ajusté. Un
\emph{hyperparamètre}\index{hyperparamètre} (hyperparameter) est un choix que l'ajustement ne fait pas lui-même (une pénalité, une profondeur d'arbre, un nombre
de couches, un taux d'apprentissage)~; l'\emph{ensemble de validation}\index{ensemble de validation} (validation set) est constitué de données mises à l'écart de l'ajustement, et sur lesquelles
les hyperparamètres sont choisis. L'\emph{ensemble de test}\index{ensemble de test} (test set) est utilisé une seule fois, après tous les choix, pour
estimer comment le modèle retenu se comportera.
\end{definition}

Sur les données de marché, les trois ensembles sont des tranches de temps consécutives, dans cet ordre (\cref{fig:ml:why:split}), parce que le
modèle sera utilisé sur un futur qu'il n'a pas vu. L'\omterm{def:ml:why-financial-machine-learning-is-different:sets}{ensemble de test} est l'échantillon réservé (holdout) du livre~7 (chapitre~1)~: il n'est touché qu'une fois. Un \omterm{def:ml:why-financial-machine-learning-is-different:sets}{ensemble de test}
consulté deux fois est devenu un \omterm{def:ml:why-financial-machine-learning-is-different:sets}{ensemble de validation}.

\begin{omfigure}
\begin{tikzpicture}[x=0.95cm,y=1cm,font=\small]
\draw[->,thick] (0,0) -- (13.2,0) node[right]{temps};
\fill[omDef!25] (0,0.2) rectangle (7,0.9);
\fill[omProp!30] (7.15,0.2) rectangle (9.5,0.9);
\fill[omThm!25] (9.65,0.2) rectangle (12.8,0.9);
\node at (3.5,0.55) {entraînement~: ajuster $\theta$};
\node at (8.33,0.55) {validation};
\node at (11.22,0.55) {test~: utilisé une fois};
\draw[omDef,thick] (3.5,-0.15) -- (3.5,-0.45) node[below,align=center]{\footnotesize variables et cibles\\[-2pt]\footnotesize connues à chaque date};
\draw[omProp,thick] (8.33,-0.15) -- (8.33,-0.45) node[below,align=center]{\footnotesize choisir les\\[-2pt]\footnotesize hyperparamètres};
\draw[omThm,thick] (11.22,-0.15) -- (11.22,-0.45) node[below,align=center]{\footnotesize estimer\\[-2pt]\footnotesize l'erreur de\\[-2pt]\footnotesize généralisation};
\end{tikzpicture}

\omcaption{Sur les données de marché, les trois ensembles sont des tranches de temps consécutives. Le chapitre~3 remplace
la découpe unique par des blocs purgés et des schémas walk-forward~; l'ordre ne change jamais.}\label{fig:ml:why:split}
\end{omfigure}

\begin{definition}[Erreur de généralisation, surapprentissage (overfitting), décomposition biais--variance]\label{def:ml:why-financial-machine-learning-is-different:generalisation}
L'\emph{erreur de généralisation}\index{erreur de généralisation} (generalisation error) d'un modèle ajusté $\hat f$ est sa perte espérée sur un nouveau
tirage $(x, y)$ dans la distribution à laquelle le modèle sera confronté. Le \emph{surapprentissage}\index{surapprentissage} (overfitting) est l'ajustement de
variables de l'échantillon d'entraînement qui ne se reproduisent pas, de sorte que la perte dans l'échantillon (in-sample) sous-estime l'erreur de généralisation.
Pour la perte quadratique, la \emph{décomposition biais--variance}\index{décomposition biais--variance} (bias--variance decomposition) découpe, en un point, l'erreur de généralisation
espérée en bruit, biais au carré et variance (\cref{prop:ml:why:bv}).
\end{definition}

\begin{proposition}[Décomposition biais--variance]\label{prop:ml:why:bv}
Soit $y = \mu(x) + \varepsilon$ avec $\E[\varepsilon\mid x] = 0$ et $\Var(\varepsilon\mid x) = \sigma^2(x)$, et soit
$\hat f$ ajusté sur un échantillon d'entraînement indépendant de $(x, y)$. Alors
\[
\E\bigl[(y - \hat f(x))^2 \bigm| x\bigr] = \sigma^2(x) + \bigl(\E[\hat f(x)] - \mu(x)\bigr)^2 + \Var\bigl(\hat f(x)\bigr),
\]
où $\E$ et $\Var$ de $\hat f(x)$ portent sur les échantillons d'entraînement.
\end{proposition}

\begin{proof}
Écrivons $y - \hat f = \varepsilon + (\mu - \E\hat f) + (\E\hat f - \hat f)$. Les trois termes sont non corrélés conditionnellement à $x$~:
$\varepsilon$ est indépendant de l'échantillon d'entraînement et d'espérance nulle, et $\E\hat f - \hat f$ est d'espérance nulle sur
les échantillons d'entraînement, tandis que $\mu - \E\hat f$ est une constante. Élever au carré et prendre l'espérance.
\end{proof}

Le bruit $\sigma^2$ est le plancher en deçà duquel aucun modèle ne descend. Dans la plupart des problèmes d'apprentissage automatique, le bruit est faible et le combat se joue
entre le biais (un modèle trop rigide pour $\mu$) et la variance (un modèle qui suit l'échantillon). Sur les données de marché, $\sigma^2$ représente
presque toute la perte, et tout modèle doté d'une variance, même faible, la consacre au bruit~: l'équilibre penche vers des modèles rigides,
régularisés (chapitre~4).

\section{Rapport signal sur bruit proche de zéro et plafond du R²}

\begin{definition}[R² hors échantillon, plafond de prévisibilité]\label{def:ml:why-financial-machine-learning-is-different:r2}
Le \emph{R² hors échantillon}\index{R2 hors échantillon} (out-of-sample R-squared) des prévisions $\hat y_i$ des rendements $y_i$ sur un
\omterm{def:ml:why-financial-machine-learning-is-different:sets}{ensemble de test} est
\[
R^2_{\mathrm{oos}} = 1 - \frac{\sum_i(y_i - \hat y_i)^2}{\sum_i y_i^2},
\]
mesuré par rapport à une prévision nulle, et non à la moyenne de l'échantillon. Le \emph{plafond de prévisibilité}\index{plafond de prévisibilité} (predictability ceiling)
d'un jeu de données est le $R^2_{\mathrm{oos}}$ de l'espérance conditionnelle vraie $\mu(x_i) = \E[y_i\mid x_i]$~: le meilleur score
qu'un modèle de ces variables puisse atteindre.
\end{definition}

L'indice de référence nul est celui de Gu, Kelly et Xiu, dont l'étude de 30\,000 actions sur 60 ans l'utilise parce que la moyenne historique
du rendement d'une action est une prévision si bruitée qu'il est trop facile de la battre. Leur R² mensuel hors échantillon, calculé action par action,
va de 0.16~\% pour un modèle linéaire à trois variables à 0.33--0.40~\% pour les arbres et les réseaux de neurones, avec
0.40~\% pour le meilleur réseau. Ce sont là les chiffres d'un domaine où un demi-pour-cent est un résultat solide.

\begin{proposition}[Le plafond borne tout modèle~; R² et corrélation]\label{prop:ml:why:ceiling}
\begin{enumerate}
\item Pour toute prévision $\hat f(x)$ construite à partir des variables, $\E(y - \hat f)^2 \ge \E(y - \mu)^2$~: en espérance, aucun
modèle ne voit son $R^2_{\mathrm{oos}}$ dépasser le plafond.
\item Si $y$ et une prévision $f$ sont d'espérance nulle, $\operatorname{Corr}(f, y) = \rho$, et si la prévision est utilisée sous la forme $c f$, alors
$R^2_{\mathrm{oos}}(c) = (2c\rho\sigma_y\sigma_f - c^2\sigma_f^2)/\sigma_y^2$, maximal en $c^\star = \rho\sigma_y/\sigma_f$
où il vaut $\rho^2$, et nul en $c = 2c^\star$.
\end{enumerate}
\end{proposition}

\begin{proof}
(1) $\E(y - \hat f)^2 = \E(y - \mu)^2 + \E(\mu - \hat f)^2$ car $y - \mu$ est non corrélé avec toute fonction de $x$.
(2) Développer $\E(y - cf)^2 = \sigma_y^2 - 2c\rho\sigma_y\sigma_f + c^2\sigma_f^2$ et maximiser le trinôme en $c$.
\end{proof}

La seconde partie relie le R² au coefficient d'information du livre~7 (chapitre~6)~: une prévision dont la corrélation
avec les rendements est $\rho$ capte au plus $\rho^2$, et une corrélation de 0.075 correspond à un R² de 0.56~\%. Elle indique aussi comment une
bonne prévision est mal notée~: une prévision au bon sens et dont la taille est le double de la bonne obtient zéro.

Le bac à sable de l'ouvrage, \texttt{firm.mlsynth}, tire une coupe transversale mensuelle de 500 actions dotées de vingt caractéristiques
publiées sous forme de rangs transversaux dans $[-1, 1]$, comme la littérature empirique les classe. Six portent un signal, de façon linéaire et
par une interaction, un carré, un seuil et une valeur absolue~; quatorze ne sont que du bruit. Les rendements ajoutent un facteur de marché,
dix facteurs sectoriels et des chocs spécifiques à queue épaisse, et le rendement espéré injecté est ajusté de sorte que le plafond
soit de 0.7~\%. La régression ridge, les arbres boostés par gradient de LightGBM et un petit perceptron multicouche sont entraînés sur vingt
ans et testés sur les dix suivantes (\cref{tab:ml:why:models})~; pour comparaison, les trois mêmes modèles sont ajustés à une tâche
avec un R² optimal au sens de Bayes de 0.95 (20\,000 exemples, une fonction lisse et non linéaire de cinq des vingt variables).

\begin{table}[htbp]
\centering\small
\setlength{\tabcolsep}{4pt}%
\begin{tabular}{@{}>{\raggedright\arraybackslash}p{3.1cm}>{\centering\arraybackslash}p{2.0cm}>{\centering\arraybackslash}p{2.0cm}>{\centering\arraybackslash}p{2.0cm}>{\centering\arraybackslash}p{2.0cm}>{\centering\arraybackslash}p{2.0cm}@{}}
\toprule
 & \multicolumn{3}{c}{\shortstack{500 actions,\\ plafond 0.60~\% hors échantillon}} & \multicolumn{2}{c}{signal fort, Bayes 0.95}\\
\cmidrule(lr){2-4}\cmidrule(l){5-6}
modèle & dans l'échantillon & hors échantillon & corr.\ avec la vérité & dans l'échantillon & hors échantillon\\
\midrule
régression ridge & 0.39~\% & 0.38~\% & 0.74 & 0.72 & 0.72\\
arbres boostés & 1.45~\% & 0.57~\% & 0.92 & 0.92 & 0.92\\
réseau de neurones & 0.64~\% & 0.28~\% & 0.74 & 0.94 & 0.93\\
\bottomrule
\end{tabular}
\caption{Trois modèles sur deux tâches~: vingt ans d'entraînement et dix ans de test sur les 500 actions de \texttt{firm.mlsynth},
et 20\,000 exemples d'entraînement d'une tâche à signal fort. R² par rapport à zéro pour les actions, par rapport à la moyenne pour
la tâche générique. Données~: \texttt{ml\_why.compare} et \texttt{ml\_why.strong\_signal}.}\label{tab:ml:why:models}
\end{table}

Le \Cref{tab:ml:why:models} est le chapitre en miniature. Sur la tâche à signal fort, les scores dans l'échantillon et hors échantillon de chaque modèle
concordent et les modèles flexibles l'emportent. Sur les actions, le R² dans l'échantillon des arbres boostés atteint 1.45~\%, soit deux fois le plafond
des années d'entraînement (0.72~\%)~: l'excédent est du bruit que le modèle a ajusté. Son \omterm{def:ml:why-financial-machine-learning-is-different:r2}{R² hors échantillon}, 0.57~\%, représente 95~\% du plafond, le
meilleur des trois~; le modèle linéaire en atteint 64~\%, et le réseau, entraîné ici sans le soin du
chapitre~7, moins de la moitié. Les corrélations des prévisions avec la vérité
(de 0.74 à 0.92) sont élevées~; leurs corrélations avec les rendements réalisés sont inférieures à 0.08.

\begin{omfigure}
\begin{tikzpicture}
\begin{axis}[width=11cm,height=5.2cm,ybar,bar width=4pt,xlabel={R² mensuel hors échantillon (\%)},ylabel={mois},
  tick label style={font=\small},label style={font=\small},xmin=-1.1,xmax=2.85,ymin=0,grid=major,grid style={gray!25},
  legend columns=2,legend style={font=\footnotesize,at={(0.5,-0.32)},anchor=north,draw=none,fill=none,
  /tikz/every even column/.append style={column sep=8pt}},table/col sep=comma]
\addplot[fill=omDef!55,draw=omDef] table[x=x,y=boosting]{figdata/ml/01-why-financial-machine-learning-is-different/months.csv};
\addplot[fill=gray!45,draw=gray] table[x=x,y=truth]{figdata/ml/01-why-financial-machine-learning-is-different/months.csv};
\draw[omThm,thick,dashed] (axis cs:0,0) -- (axis cs:0,26);
\legend{arbres boostés,la vérité $\mu$}
\end{axis}
\end{tikzpicture}

\omcaption{Les 120 mois de test, un R² par mois sur 500 actions, pour les arbres boostés et pour la vérité injectée
elle-même. Les barres à gauche de la ligne pointillée sont les mois où la prévision a fait moins bien que zéro~: 18~\% d'entre eux pour le modèle, et
quelques-uns pour la vérité. Données~: \texttt{ml\_why.compare} sur \texttt{firm.mlsynth}.}\label{fig:ml:why:months}
\end{omfigure}

\section{Non-stationnarité et rétroaction adverse}

La distribution qu'un modèle rencontre n'est pas celle sur laquelle il a été entraîné. Certains changements sont exogènes (une nouvelle réglementation, une nouvelle
plateforme de négociation, une crise). D'autres sont causés par les transactions du modèle et celles de ses concurrents~: un prédicteur public, ou que
plusieurs entreprises ont trouvé, est exploité jusqu'à ce que son rendement devienne inférieur à son coût. Le livre~7 a mesuré ce résultat sur des anomalies réelles,
sous la forme d'une décroissance après publication (chapitre~13), et sa cause, l'encombrement (crowding, chapitre~28). Dans les sciences physiques, les
données ne lisent pas l'article.

\texttt{firm.mlsynth} peut faire vieillir un signal. Dès le premier mois de test, l'effet de la caractéristique linéaire la plus forte décroît
avec une demi-vie de deux ans, et cinq ans après le début du test, l'interaction entre deux caractéristiques change de signe.
Les arbres boostés, entraînés une seule fois sur les vingt premières années, conservent leur R² malgré la décroissance, que
masquent les autres caractéristiques, et le perdent lorsque l'interaction change de signe (\cref{fig:ml:why:drift})~: les années 6 à 10 affichent en moyenne
0.01~\%, tandis que le plafond, qui se déplace avec la nouvelle vérité, vaut en moyenne 0.83~\%. Rien dans les entrées du modèle n'a annoncé
le changement~; ses erreurs, elles, l'ont fait. Le chapitre~12 le détecte et
le chapitre~27 construit les alarmes.

\begin{omfigure}
\begin{tikzpicture}
\begin{axis}[width=11cm,height=5.4cm,xlabel={année de test},ylabel={R² hors échantillon (\%)},
  tick label style={font=\small},label style={font=\small},xmin=0.5,xmax=10.5,ymin=-0.3,ymax=1.1,
  xtick={1,2,3,4,5,6,7,8,9,10},grid=major,grid style={gray!25},
  legend columns=3,legend style={font=\footnotesize,at={(0.5,-0.3)},anchor=north,draw=none,fill=none,
  /tikz/every even column/.append style={column sep=8pt}},table/col sep=comma]
\addplot[omDef,very thick,mark=*] table[x=year,y=stationary]{figdata/ml/01-why-financial-machine-learning-is-different/drift.csv};
\addplot[omThm,very thick,mark=square*] table[x=year,y=changing]{figdata/ml/01-why-financial-machine-learning-is-different/drift.csv};
\addplot[black!60,thick,dashed] table[x=year,y=ceiling_changing]{figdata/ml/01-why-financial-machine-learning-is-different/drift.csv};
\draw[gray,thin] (axis cs:0.5,0) -- (axis cs:10.5,0);
\legend{marché stationnaire,signal décroît puis s'inverse,plafond (marché changeant)}
\end{axis}
\end{tikzpicture}

\omcaption{Arbres boostés entraînés sur les années 1--20 de deux panels avec la même graine, évalués année par année sur les dix suivantes.
Dans le second, l'effet linéaire de la caractéristique~0 décroît à partir de l'année de test~1 (demi-vie de deux ans) et l'interaction des
caractéristiques 0 et 3 change de signe au début de l'année~6. Données~: \texttt{ml\_why.drift\_by\_year}.}\label{fig:ml:why:drift}
\end{omfigure}

\section{Petits échantillons effectifs}

L'\omterm{def:ml:why-financial-machine-learning-is-different:sets}{ensemble d'entraînement} compte $240 \times 500 = 120\,000$ mois-actions. Ce ne sont pas 120\,000 observations indépendantes. Les
rendements imprévus des actions partagent un marché et un secteur~: leur corrélation moyenne deux à deux est de 0.27, et le
nombre effectif d'actions indépendantes dans un mois, $n/(1 + (n-1)\bar\rho)$, vaut 3.7 sur 500. Tout ce qui dépend de la
trajectoire du marché est estimé à partir de 240 tirages, et non de 120\,000.

L'ordonnée à l'origine est le cas le plus clair. Entraîné sur des rendements bruts, l'ordonnée à l'origine d'un modèle est la moyenne du rendement de marché
sur la fenêtre d'entraînement. La prime vraie dans le bac à sable est nulle~; la fenêtre a affiché en moyenne $-0.56~\%$ par mois, soit environ deux
erreurs types d'une moyenne sur 240 mois, à une volatilité mensuelle de 4.5~\%. Les trois modèles entraînés sur des rendements bruts obtiennent $-0.29~\%$,
$-0.09~\%$ et $-0.30~\%$ hors échantillon, chacun moins bien qu'une prévision nulle. Entraînés sur les rendements moins la
moyenne transversale de chaque mois, ils apprennent quelles actions font mieux, non où va le marché, et obtiennent les chiffres du tableau.

\begin{method}[Avant d'ajuster un modèle de rendements]\label{met:ml:why:before}
\begin{enumerate}
\item Calculer le plafond, ou en borner la valeur~: sur des données synthétiques, exactement~; sur des données réelles, à partir du meilleur R² publié ou
du coefficient d'information de travaux comparables (\cref{prop:ml:why:ceiling}).
\item Décider ce qu'on demande au modèle~: la coupe transversale (centrer ou classer la cible par date), ou le niveau.
\item Compter l'échantillon effectif~: des dates et non des lignes pour tout ce qui est commun à tous les actifs~; les étiquettes qui se chevauchent
(chapitre~2) la divisent à nouveau.
\item Fixer la période de test avant d'en regarder les résultats, et prévoir la durée d'un suivi en conditions réelles nécessaire pour distinguer le modèle de zéro et
du modèle plus simple.
\end{enumerate}
\end{method}

La durée est l'autre échantillon effectif. Le R² mensuel des arbres boostés a une moyenne de 0.62~\% et un écart type
de 0.64~\%~: 4.3 mois de négociation en conditions réelles placent sa moyenne à deux erreurs types de zéro. Distinguer ce modèle de la régression ridge demande
davantage de temps~: la différence mensuelle a une moyenne de 0.21~\% et un écart type de 0.47~\%, ce qui exige 21 mois. Avec un historique plus court, le modèle
flexible perd son avantage (la \cref{fig:ml:why:learning})~: entraînés sur les seules deux dernières années, les arbres
obtiennent 0.18~\% et la régression ridge 0.29~\%~; les deux courbes se croisent entre quatre et huit ans. Plus le modèle est flexible, plus une part
de ses données d'entraînement sert à apprendre ce que le bruit n'est pas.

\begin{omfigure}
\begin{tikzpicture}
\begin{axis}[width=11cm,height=5.4cm,xlabel={fenêtre d'entraînement (mois)},ylabel={R² hors échantillon (\%)},
  tick label style={font=\small},label style={font=\small},xmin=0,xmax=252,ymin=0,ymax=0.7,xtick={24,48,96,144,240},
  grid=major,grid style={gray!25},legend columns=3,legend style={font=\footnotesize,at={(0.5,-0.3)},anchor=north,
  draw=none,fill=none,/tikz/every even column/.append style={column sep=8pt}},table/col sep=comma]
\addplot[omDef,very thick,mark=*] table[x=months,y=boosting]{figdata/ml/01-why-financial-machine-learning-is-different/learning.csv};
\addplot[omThm,very thick,mark=square*] table[x=months,y=ridge]{figdata/ml/01-why-financial-machine-learning-is-different/learning.csv};
\addplot[black!60,thick,dashed] table[x=months,y=ceiling]{figdata/ml/01-why-financial-machine-learning-is-different/learning.csv};
\legend{arbres boostés,régression ridge,plafond}
\end{axis}
\end{tikzpicture}

\omcaption{Courbes d'apprentissage sur le panel d'actions~: les mêmes dix années de test, avec des fenêtres d'entraînement
allant des 2 aux 20 dernières années. Données~: \texttt{ml\_why.learning\_curve}.}\label{fig:ml:why:learning}
\end{omfigure}

\section{Ce qui se transporte du reste de l'apprentissage automatique, et ce qui ne se transporte pas}

Les algorithmes se transportent~: les moindres carrés et leurs pénalisations (livre~4, chapitre~16), les arbres, les réseaux de neurones, les méthodes de gradient
(livre~4, chapitre~24). Il en va de même de la discipline de l'évaluation sur données réservées. Ce qui ne se transporte pas, ce sont les réglages par défaut
et les intuitions forgées sur les images et le texte~: la validation croisée avec mélange (livre~7, chapitre~20, et
chapitre~3), un score de validation lu comme un fait plutôt que comme un tirage, des \omterm{def:ml:why-financial-machine-learning-is-different:sets}{hyperparamètres} ajustés jusqu'à ce que
l'\omterm{def:ml:why-financial-machine-learning-is-different:sets}{ensemble de validation} donne raison, une exactitude de 55~\% prise pour un modèle faible alors qu'elle peut être proche du plafond, et un modèle ajusté
une fois puis gardé pendant des années. Chaque chapitre de ce livre est une de ces corrections, mesurée.

\section{Tutoriel~: le modèle à 0.57~\%}\label{tut:ml:why-financial-machine-learning-is-different:tut}

\begin{tutorial}
\textbf{Objectif.} Mesurer le plafond d'un panel d'actions synthétique, comparer trois modèles dans et hors échantillon, et refaire le calcul sur
une tâche à fort signal. \textbf{État final~:} le tableau, les \cref{fig:ml:why:months,fig:ml:why:drift,fig:ml:why:learning}.
\begin{tutsteps}
\item \textbf{Le panel.} Les caractéristiques sont des rangs de processus latents persistants~; la vérité est une part linéaire et une
part non linéaire, mises à l'échelle pour que $R^2_{\mathrm{oos}}(r,\mu)$ atteigne le plafond.
\omcode{code/firm/mlsynth/firm_mlsynth.py}{131}{163}{Le panel d'actions synthétique avec un rendement espéré incorporé.}\label{lst:ml:why:panel}
\item \textbf{La cible et les trois modèles.} Entraîner sur les rendements diminués de la moyenne transversale de chaque mois~; évaluer par rapport à
zéro.
\omcode{code/ml/01-why-financial-machine-learning-is-different/python/ml_why.py}{75}{91}{R² dans et hors échantillon, et plafond.}\label{lst:ml:why:compare}
\item \textbf{Exécuter} \texttt{compare()}, \texttt{compare(target='raw')}, \texttt{strong\_signal()},
\texttt{learning\_curve()}, \texttt{effective\_stocks()}, \texttt{drift\_by\_year()} et \texttt{fig\_why.py}.
\end{tutsteps}
\textbf{Pour aller plus loin.} Abaisser le rapport signal sur bruit de la tâche forte jusqu'à ce que le R² de Bayes atteigne 1~\% et observer
quel modèle l'emporte avec 20\,000 exemples (exercice~7)~; porter
\texttt{nonlin} à 0.9 et voir si la régression ridge atteint encore la moitié du plafond.
\end{tutorial}

\section{Construction~: tâches d'apprentissage synthétiques}\label{bld:ml:why-financial-machine-learning-is-different:mlsynth}

\begin{build}
\bfield{Rôle} Chaque modèle de ce livre est évalué par rapport au meilleur prédicteur qui existe pour ses données.
\bfield{Interface} \texttt{task(n, p, snr, kind, seed)} avec son R² de Bayes~;
\texttt{PanelConfig(n, months, k, seed, ceiling, nonlin, premium, mkt\_vol, ind\_vol, spec\_vol, t\_df, decay\_from, decay\_half\_life, regime\_at)}~;
\texttt{panel(cfg) -> Panel} avec \texttt{X} (mois, actions, caractéristiques), \texttt{r}, \texttt{mu},
\texttt{beta}, \texttt{industry}, \texttt{mkt}, et \texttt{Panel.flat(months)}~; \texttt{r2\_oos},
\texttt{ceiling}, \texttt{months\_to\_detect}.
\bfield{Règles} Les caractéristiques au mois $t$ sont connues à sa fin et prédisent le rendement du mois $t+1$~; la vérité n'est
jamais une entrée~; graine fixée, calcul déterministe.
\bfield{Tests d'acceptation} \texttt{code/firm/mlsynth/tests/}~: la vérité atteint le R² de Bayes sur la tâche générique~; le
plafond est atteint en moyenne sur les graines~; les caractéristiques de bruit n'apportent (presque) rien~; la décroissance et le retournement modifient la
vérité~; les formules du R² et de la détection sur de petits cas.
\bfield{Pour aller plus loin} Caractéristiques observées avec erreur, de sorte que le plafond obtenu à partir des variables soit inférieur à celui de la vérité~;
valeurs manquantes et cotations qui commencent et s'arrêtent, comme dans \texttt{firm.synthmkt}.
\end{build}

\begin{omsources}
\item S.~Gu, B.~Kelly et D.~Xiu, ``Empirical asset pricing via machine learning'', \emph{Review of Financial Studies}
33(5), 2020 (NBER Working Paper 25398, 2018, révisé en 2019, tableau~1).
\item S.~Geman, E.~Bienenstock et R.~Doursat, ``Neural networks and the bias/variance dilemma'', \emph{Neural
Computation} 4(1), 1992.
\item T.~Hastie, R.~Tibshirani et J.~Friedman, \emph{The Elements of Statistical Learning}, 2\up{e} éd., Springer, 2009.
\item R.~Israel, B.~Kelly et T.~Moskowitz, ``Can machines `learn' finance?'', \emph{Journal of Investment Management},
2020 (SSRN 3624052).
\item M.~López de Prado, \emph{Advances in Financial Machine Learning}, Wiley, 2018.
\end{omsources}

\section{Exercices}

\begin{exercise}[$\star$]\label{exo:ml:why-financial-machine-learning-is-different:1}
Une prévision des rendements du mois suivant a une corrélation de 0.05 avec eux et est optimalement mise à l'échelle. Quel est son
\omterm{def:ml:why-financial-machine-learning-is-different:r2}{R² hors échantillon}~? Qu'en est-il si on l'utilise à trois fois l'échelle optimale~?
\end{exercise}

\begin{exercise}[$\star$]\label{exo:ml:why-financial-machine-learning-is-different:2}
Les rendements non anticipés de 500 actions ont une corrélation moyenne deux à deux de 0.27. Quel est le nombre effectif
d'actions indépendantes sur un mois~? Et avec 0.05~?
\end{exercise}

\begin{exercise}[$\star$]\label{exo:ml:why-financial-machine-learning-is-different:3}
Le R² mensuel d'un modèle a une moyenne de 0.5~\% et un écart type de 0.8~\%. Combien de mois de résultats indépendants
faut-il pour placer la moyenne à deux erreurs types de zéro~? Trois~?
\end{exercise}

\begin{exercise}[$\star\star$]\label{exo:ml:why-financial-machine-learning-is-different:4}
Les arbres boostés obtiennent 1.45~\% dans l'échantillon, alors que le plafond est de 0.72~\%. Expliquer comment un modèle peut battre le meilleur prédicteur possible
dans l'échantillon, et ce que l'écart indique sur son score hors échantillon.
\end{exercise}

\begin{exercise}[$\star\star$]\label{exo:ml:why-financial-machine-learning-is-different:5}
Le rendement moyen du marché sur la fenêtre d'entraînement était de $-0.56~\%$ par mois, et l'espérance du carré du rendement mensuel d'une action est
$\E r^2 = 0.0093$. Quelle est l'erreur type d'une moyenne sur 240 mois pour une volatilité mensuelle
de 4.5~\%, et combien de points de R² une erreur d'ordonnée à l'origine de $-0.56~\%$ coûte-t-elle~?
\end{exercise}

\begin{exercise}[$\star\star$]\label{exo:ml:why-financial-machine-learning-is-different:6}
\emph{Trouver la faille.} «~Nous avons entraîné notre modèle sur des rendements mensuels bruts, et son \omterm{def:ml:why-financial-machine-learning-is-different:r2}{R² hors échantillon} est négatif, donc les caractéristiques
ne portent aucune information sur les rendements.~»
\end{exercise}

\begin{exercise}[$\star\star\star$]\label{exo:ml:why-financial-machine-learning-is-different:7}
\emph{Programmation.} Lancer \texttt{ml\_why.strong\_signal} avec un rapport signal sur bruit qui donne un R² de Bayes de 1~\%
(\texttt{snr = 0.0101/0.9899}), 20\,000 exemples d'entraînement. Rapporter le \omterm{def:ml:why-financial-machine-learning-is-different:r2}{R² hors échantillon} des trois modèles et expliquer
le classement.
\end{exercise}

\begin{exercise}[$\star\star\star$]\label{exo:ml:why-financial-machine-learning-is-different:8}
Démontrer la partie 2 de la \cref{prop:ml:why:ceiling} et s'en servir pour expliquer pourquoi un modèle dont les prévisions sont corrélées à 0.92 avec
la vérité peut néanmoins obtenir un score négatif hors échantillon.
\end{exercise}

\section{Problème~: Le modèle à 0.57~\%}

\begin{problem}[{Problème du week-end --- un modèle proche de son plafond}]\label{pb:ml:why-financial-machine-learning-is-different:1}
Le panel du chapitre~: 500 actions, vingt caractéristiques, vingt ans d'entraînement et dix ans de test, ainsi qu'une vérité introduite à dessein.

\medskip\noindent\textbf{Partie I --- Le plafond.}
\begin{enumerate}
\item Quel est le \omterm{def:ml:why-financial-machine-learning-is-different:r2}{plafond de prévisibilité} sur les années d'entraînement et sur celles de test, et pourquoi diffèrent-ils~?
\item À quelle corrélation avec les rendements correspond un R² de 0.60~\%~?
\item Pourquoi mesure-t-on le R² par rapport à zéro plutôt qu'à la moyenne de l'échantillon~?
\item Quelle part des mois la vérité elle-même perd-elle face à une prévision nulle, d'après la \cref{fig:ml:why:months}~?
\end{enumerate}

\medskip\noindent\textbf{Partie II --- Les modèles.}
\begin{enumerate}[resume]
\item Donner le R² dans et hors échantillon des trois modèles.
\item Quelle part du plafond chaque modèle atteint-il hors échantillon~?
\item Pourquoi le classement sur la tâche forte diffère-t-il du classement sur les actions~?
\item Quel terme de la \omterm{def:ml:why-financial-machine-learning-is-different:generalisation}{décomposition biais--variance} explique le 1.45~\% dans l'échantillon des arbres~?
\end{enumerate}

\medskip\noindent\textbf{Partie III --- Échantillons effectifs.}
\begin{enumerate}[resume]
\item Combien d'actions indépendantes vaut un mois, et pourquoi~?
\item Quels scores obtiennent les modèles entraînés sur les rendements bruts, et pourquoi~?
\item Quelle durée de fenêtre d'entraînement le boosting doit-il avoir pour battre la régression ridge~?
\item Combien de mois de résultats en production faut-il pour distinguer les arbres de zéro et de la ridge~?
\end{enumerate}

\medskip\noindent\textbf{Partie IV --- Le verdict.}
\begin{enumerate}[resume]
\item Que se passe-t-il pour les arbres lorsque l'interaction change de signe, et que montraient les entrées~?
\item Énoncer le \emph{résultat nommé}~: le plafond, la part que le modèle en atteint, et les mois nécessaires pour le distinguer de zéro
et de la ridge.
\item Faut-il trader les arbres ou la régression ridge, avec dix ans d'historique~? Avec deux~?
\item Que devrait surveiller une règle de suivi, compte tenu de la première question de la partie IV~?
\item Comment l'encombrement (crowding) apparaîtrait-il dans le bac à sable~?
\item Pourquoi le plafond est-il inconnaissable sur des données réelles, et par quoi le remplace-t-on~?
\item Que suggère le meilleur R² mensuel de 0.40~\% de Gu, Kelly et Xiu sur les plafonds réels~?
\item En une phrase~: pourquoi l'apprentissage automatique financier est-il différent~?
\end{enumerate}
\end{problem}

\section{Questions d'entretien}

\begin{interviewq}[$\star$ \iqroles{chercheur, ingénieur ML}]\label{iq:ml:why-financial-machine-learning-is-different:1}
Votre modèle de rendements mensuels d'actions a un \omterm{def:ml:why-financial-machine-learning-is-different:r2}{R² hors échantillon} de 0.5~\%. Est-ce bon~?
\end{interviewq}

\begin{interviewq}[$\star$ \iqroles{chercheur, ingénieur ML}]\label{iq:ml:why-financial-machine-learning-is-different:2}
Énoncez la \omterm{def:ml:why-financial-machine-learning-is-different:generalisation}{décomposition biais--variance}. Où les modèles de rendements se situent-ils sur ce compromis, et pourquoi~?
\end{interviewq}

\begin{interviewq}[$\star\star$ \iqroles{chercheur}]\label{iq:ml:why-financial-machine-learning-is-different:3}
Vous disposez de dix millions de lignes de données quotidiennes sur 2\,000 actions sur 20 ans. Combien d'observations indépendantes avez-vous~?
\end{interviewq}

\begin{interviewq}[$\star\star$ \iqroles{chercheur}]\label{iq:ml:why-financial-machine-learning-is-different:4}
Pourquoi les études de prédiction des rendements mesurent-elles le R² par rapport à une prévision nulle plutôt qu'à la moyenne historique~?
\end{interviewq}

\begin{interviewq}[$\star\star$ \iqroles{chercheur, trader}]\label{iq:ml:why-financial-machine-learning-is-different:5}
Un modèle avait un coefficient d'information de 0.06 dans son backtest et de 0.02 pendant ses six premiers mois en production. Énumérez les
causes possibles et dire comment vous les distingueriez.
\end{interviewq}

\begin{interviewq}[$\star\star\star$ \iqroles{chercheur, ingénieur ML}]\label{iq:ml:why-financial-machine-learning-is-different:6}
Montrez qu'une prévision de signe correct peut avoir un \omterm{def:ml:why-financial-machine-learning-is-different:r2}{R² hors échantillon} négatif, et dites ce que vous feriez pour y remédier.
\end{interviewq}
